% AA-06 | Algebra abstrakcyjna | Iteracja 03
\chapter[Iloczyny i rozszerzenia grup]{Iloczyny proste, półproste i rozszerzenia grup}
\label{aa:ch:konstrukcje-grup}

\section{Iloczyn prosty zewnętrzny}
\begin{definicja}[Iloczyn prosty grup]
\label{aa:def:directproduct}
Niech $G,H$ będą grupami o elementach neutralnych $e_G,e_H$.
Na zbiorze kartezjańskim $G\times H$ definiujemy mnożenie
$(g,h)(g',h')=(gg',hh')$.
Tak otrzymany obiekt nazywamy \emph{zewnętrznym iloczynem prostym}
$G\times H$.
\end{definicja}
\begin{twierdzenie}[Aksjomaty i struktura iloczynu prostego]
\label{aa:tw:direct-product}
Dla dowolnych grup $G,H$ iloczyn $G\times H$ jest grupą,
której jedynką jest $(e_G,e_H)$, a odwrotnością
$(g,h)^{-1}=(g^{-1},h^{-1})$.
Odwzorowania $\pi_G(g,h)=g$, $\pi_H(g,h)=h$ są surjektywnymi
homomorfizmami. Podgrupy $G\times\{e_H\}$ i
$\{e_G\}\times H$ są normalne, mają trywialne przecięcie
i komutują element po elemencie.
Ponadto $Z(G\times H)=Z(G)\times Z(H)$.
\end{twierdzenie}
\begin{proof}
Łączność wynika z równości
\begin{align*}
((g,h)(g',h'))(g'',h'')&=((gg')g'',(hh')h''),\\
(g,h)((g',h')(g'',h''))&=(g(g'g''),h(h'h'')).
\end{align*}
Obie pary są równe z łączności mnożenia w $G$ i $H$.
Mnożenie przez $(e_G,e_H)$ pozostawia każdą parę bez zmiany,
a podana odwrotność jest obustronna, gdyż w każdej współrzędnej
mamy iloczyn elementu i jego odwrotności.
Projekcje zachowują mnożenie z definicji. Są surjektywne,
ponieważ $\pi_G(g,e_H)=g$ dla każdego $g\in G$
oraz $\pi_H(e_G,h)=h$ dla każdego $h\in H$.
Jądro $\pi_H$ to $G\times\{e_H\}$,
a jądro $\pi_G$ to $\{e_G\}\times H$, więc podgrupy
te są normalne na mocy twierdzenia o jądrze homomorfizmu.
Ich przekrój to jedynka. Równości
$(g,e_H)(e_G,h)=(g,h)=(e_G,h)(g,e_H)$
pokazują przemienność obu czynników.
Wreszcie $(g,h)$ komutuje z każdą $(x,y)$ dokładnie wtedy,
gdy $gx=xg$ dla każdego $x\in G$ oraz $hy=yh$ dla każdego
$y\in H$. To oznacza $g\in Z(G)$ i $h\in Z(H)$.
\end{proof}

\begin{stwierdzenie}[Własność odwzorowań do iloczynu]
\label{aa:st:product-universal}
Dla grup $K,G,H$ oraz homomorfizmów $f:K\to G$,
$g:K\to H$ istnieje dokładnie jeden homomorfizm
$u:K\to G\times H$ spełniający $\pi_G\circ u=f$
i $\pi_H\circ u=g$, mianowicie $u(k)=(f(k),g(k))$.
\end{stwierdzenie}
\begin{proof}
Dla $a,b\in K$ mamy
$u(ab)=(f(a)f(b),g(a)g(b))=u(a)u(b)$.
Jeśli $v$ spełnia warunki, jego obie współrzędne są wymuszone:
$v(k)=(\pi_G(v(k)),\pi_H(v(k)))=(f(k),g(k))=u(k)$.
\end{proof}

\begin{przyklad}[Dlaczego $C_6\cong C_2\times C_3$, ale $C_4\not\cong C_2\times C_2$?]
Niech $C_n=\mathbb Z/n\mathbb Z$ z dodawaniem.
Odwzorowanie $f:C_6\to C_2\times C_3$,
$f([k]_6)=([k]_2,[k]_3)$, jest dobrze określone,
bo przystawanie modulo $6$ pociąga przystawanie
modulo $2$ i modulo $3$; ponadto jest homomorfizmem.
Jeśli $f([k]_6)=(0,0)$, to $2\mid k$ i $3\mid k$,
więc $6\mid k$. Jądro jest trywialne, toteż $f$ jest
różnowartościowy. Obie grupy mają sześć elementów,
zatem $f$ jest izomorfizmem.
Element $(1,1)$ ma rząd $6$.
Natomiast każdy element $C_2\times C_2$ spełnia
$2(x,y)=(0,0)$, więc ma rząd co najwyżej $2$.
W $C_4$ istnieje element $[1]_4$ rzędu $4$, stąd
$C_4\not\cong C_2\times C_2$.
To szczególny przypadek chińskiego twierdzenia o resztach,
które udowodnimy przy teorii pierścieni.
\end{przyklad}

\section{Iloczyn prosty wewnętrzny: kiedy grupa rozkłada się na czynniki?}
\begin{lemat}[Iloczyn podgrup i jego liczność]
\label{aa:lem:product-subgroups}
Niech $N,H\le G$, przy czym $N\trianglelefteq G$.
Wtedy $NH=\{nh:n\in N,\ h\in H\}$ jest podgrupą $G$.
Jeśli obie podgrupy są skończone, to
$|NH|=|N|\,|H|/|N\cap H|$.
\end{lemat}
\begin{proof}
Iloczyn dwóch elementów $nh,n'h'$ zapisujemy jako
$n(hn'h^{-1})(hh')$. Ponieważ $N$ jest normalna,
$hn'h^{-1}\in N$, a więc wynik należy do $NH$.
Odwrotność $(nh)^{-1}=h^{-1}n^{-1}$ możemy zapisać jako
$(h^{-1}n^{-1}h)h^{-1}\in NH$. Jedynka też należy do $NH$;
stąd kryterium podgrupy daje $NH\le G$.

Policzmy włókna odwzorowania zbiorów $N\times H\to NH$,
$(n,h)\mapsto nh$. Dwie pary mają ten sam obraz wtedy i tylko wtedy,
gdy $n^{-1}n'=hh'^{-1}\in N\cap H$.
Dla dowolnej pary $(n,h)$ i $k\in N\cap H$ para
$(nk,k^{-1}h)$ ma ten sam iloczyn $nh$.
Odwrotnie, powyższa równość pokazuje, że każda para
z tym samym obrazem jest tej postaci dla dokładnie jednego $k$.
Zatem każde włókno ma dokładnie $|N\cap H|$ elementów.
Policzenie par daje $|N|\,|H|=|NH|\,|N\cap H|$.
\end{proof}

\begin{twierdzenie}[Kryterium wewnętrznego iloczynu prostego]
\label{aa:tw:internal-direct}
Niech $G$ będzie grupą, a $N,H\trianglelefteq G$
normalnymi podgrupami takimi, że $N\cap H=\{e\}$
oraz $NH=G$. Wtedy odwzorowanie
$\Phi:N\times H\to G$, $\Phi(n,h)=nh$, jest izomorfizmem.
Odwrotnie, obrazy obu kanonicznych czynników w dowolnym
izomorfizmie $N\times H\cong G$ spełniają te warunki.
\end{twierdzenie}
\begin{proof}
Najpierw wykażemy, że elementy $N$ komutują z elementami $H$.
Dla $n\in N$, $h\in H$ komutator
$[n,h]=nhn^{-1}h^{-1}$ należy do $H$, ponieważ
$nhn^{-1}\in H$ z normalności $H$.
Należy też do $N$, gdyż $hn^{-1}h^{-1}\in N$ z normalności $N$,
więc $[n,h]=n(hn^{-1}h^{-1})\in N$.
Stąd $[n,h]\in N\cap H=\{e\}$, czyli $nh=hn$.
Dla $(n,h),(n',h')\in N\times H$ mamy zatem
$\Phi((n,h)(n',h'))=nn'hh'=nhn'h'=
\Phi(n,h)\Phi(n',h')$. Mapa jest homomorfizmem.
Jest surjektywna, gdyż $NH=G$.
Jej jądro to pary spełniające $nh=e$, czyli $n=h^{-1}\in N\cap H$.
Z założenia $n=h=e$, więc jądro jest trywialne:
$\Phi$ jest izomorfizmem.
Kierunek odwrotny wynika z własności kanonicznych podgrup
iloczynu prostego udowodnionych w Twierdzeniu~\ref{aa:tw:direct-product};
izomorfizm zachowuje normalność, przecięcia i generowanie grupy.
\end{proof}

\begin{uwaga}[Dlaczego normalność obu czynników jest istotna?]
Jeżeli tylko $N\trianglelefteq G$, a $H$ jest dowolną podgrupą
z $NH=G$ i $N\cap H=\{e\}$, to czynniki nie muszą komutować.
W takiej sytuacji otrzymamy \emph{iloczyn półprosty}.
Przykładem jest $S_3$ z normalną podgrupą $A_3$ rzędu $3$
i podgrupą $\langle(1\,2)\rangle$ rzędu $2$.
\end{uwaga}

\section{Iloczyn półprosty zewnętrzny}
\begin{definicja}[Działanie przez automorfizmy i iloczyn półprosty]
\label{aa:def:semidirect}
Niech $N,H$ będą grupami i niech
$\varphi:H\to\operatorname{Aut}(N)$ będzie homomorfizmem,
gdzie $\operatorname{Aut}(N)$ oznacza grupę automorfizmów $N$
ze składaniem odwzorowań. Dla $h\in H$ piszemy
$\varphi_h=\varphi(h)$.
Na zbiorze $N\times H$ definiujemy
\[
 (n,h)\ast(n',h')=(n\,\varphi_h(n'),\ hh').
\]
Grupę tę oznaczamy $N\rtimes_\varphi H$
i nazywamy \emph{iloczynem półprostym}.
\end{definicja}
\begin{twierdzenie}[Iloczyn półprosty jest grupą]
\label{aa:tw:semidirect-group}
Powyższe działanie jest łączne, ma jedynkę $(e_N,e_H)$,
a odwrotność pary $(n,h)$ wynosi
$(\varphi_{h^{-1}}(n^{-1}),h^{-1})$.
Podzbiór $N\times\{e_H\}$ jest normalną podgrupą izomorficzną z $N$,
a $\{e_N\}\times H$ jest podgrupą izomorficzną z $H$.
Ich przekrój jest trywialny, a razem generują całą grupę.
Jeżeli działanie $\varphi$ jest trywialne, to otrzymujemy iloczyn prosty.
\end{twierdzenie}
\begin{proof}
Porównajmy oba możliwe nawiasowania trzech elementów.
Mamy
\begin{align*}
((n,h)\ast(n',h'))\ast(n'',h'')
 &= (n\varphi_h(n')\varphi_{hh'}(n''),\ hh'h''),\\
(n,h)\ast((n',h')\ast(n'',h''))
 &= (n\varphi_h(n'\varphi_{h'}(n'')),\ hh'h'')\\
 &= (n\varphi_h(n')\varphi_h(\varphi_{h'}(n'')),\ hh'h'').
\end{align*}
Ponieważ $\varphi$ jest homomorfizmem,
$\varphi_{hh'}=\varphi_h\circ\varphi_{h'}$,
więc wyrażenia są równe. To dowodzi łączności.
Mamy $\varphi_{e_H}=\operatorname{id}_N$;
ponadto każdy automorfizm przenosi jedynkę na jedynkę.
Stąd $(e_N,e_H)$ jest jedynką dwustronną.
Dla wskazanego elementu odwrotnego obliczamy
\[
 (n,h)\ast(\varphi_{h^{-1}}(n^{-1}),h^{-1})
 =(n\varphi_h(\varphi_{h^{-1}}(n^{-1})),e_H)=(e_N,e_H).
\]
W drugą stronę pierwszy czynnik wynosi
$\varphi_{h^{-1}}(n^{-1})\varphi_{h^{-1}}(n)=e_N$,
więc otrzymujemy również jedynkę.

Włączenie $n\mapsto(n,e_H)$ jest homomorfizmem,
podobnie $h\mapsto(e_N,h)$.
Normalność pierwszego czynnika wynika z obliczenia
\[
 (n,h)\ast(n',e_H)\ast(n,h)^{-1}
 =\bigl(n\varphi_h(n')n^{-1},e_H\bigr).
\]
Wprost widać trywialność przecięcia czynników i równość
$(n,e_H)\ast(e_N,h)=(n,h)$.
Przy trywialnym $\varphi_h=\operatorname{id}_N$
wzór na mnożenie staje się współrzędnościowy,
czyli jest dokładnie mnożeniem w $N\times H$.
\end{proof}

\begin{uwaga}[Znaczenie działania]
Iloczyn półprosty pamięta nie tylko parę grup $N,H$,
lecz także homomorfizm $\varphi:H\to\operatorname{Aut}(N)$.
Różne działania mogą dawać różne grupy.
Relacja
$(e_N,h)(n,e_H)(e_N,h)^{-1}=(\varphi_h(n),e_H)$
pokazuje, że działanie $H$ na $N$ jest dokładnie
sprzężeniem przez drugi czynnik.
W odróżnieniu od iloczynu prostego dwa czynniki nie muszą komutować.
\end{uwaga}

\section{Wewnętrzny iloczyn półprosty i rozszczepione rozszerzenia}
\begin{twierdzenie}[Kryterium wewnętrznego iloczynu półprostego]
\label{aa:tw:internal-semidirect}
Niech $N\trianglelefteq G$ i $H\le G$ spełniają
$NH=G$ oraz $N\cap H=\{e\}$.
Niech $\varphi:H\to\operatorname{Aut}(N)$ będzie działaniem
przez sprzężenie, $\varphi_h(n)=hnh^{-1}$.
Wtedy $G\cong N\rtimes_\varphi H$ przez odwzorowanie
$(n,h)\mapsto nh$.
\end{twierdzenie}
\begin{proof}
Ponieważ $N$ jest normalna, sprzężenie przez element $h\in H$
ogranicza się do automorfizmu $N$.
Z równości $(hh')n(hh')^{-1}=h(h'nh'^{-1})h^{-1}$
wynika $\varphi_{hh'}=\varphi_h\circ\varphi_{h'}$.
Dla par obliczamy
\[
 (nh)(n'h')=n(hn'h^{-1})hh'=
 n\varphi_h(n')hh',
\]
co jest obrazem ich iloczynu półprostego. Mapa jest więc
homomorfizmem. Jest surjektywna z $NH=G$.
Jeśli $nh=e$, mamy $n=h^{-1}\in N\cap H=\{e\}$,
więc jądro jest trywialne; jest to izomorfizm.
Jednoznaczność przedstawienia $g=nh$ wynika też z równości
$n_1h_1=n_2h_2\implies n_2^{-1}n_1=h_2h_1^{-1}\in N\cap H$.
\end{proof}

\begin{definicja}[Krótki ciąg dokładny i rozszczepienie]
\label{aa:def:split-exact}
\emph{Krótkim ciągiem dokładnym grup} nazywamy układ
$1\longrightarrow N\xrightarrow{i}G\xrightarrow{p}Q\longrightarrow1$,
w którym $i$ jest iniektywnym homomorfizmem, $p$ surjektywnym
homomorfizmem i $\operatorname{im}i=\ker p$.
Ciąg \emph{rozszczepia się}, jeżeli istnieje homomorfizm
$s:Q\to G$ z $p\circ s=\operatorname{id}_Q$.
Takie $s$ nazywamy \emph{przekrojem homomorficznym}.
\end{definicja}

\begin{twierdzenie}[Rozszczepienie a iloczyn półprosty]
\label{aa:tw:split-semidirect}
Dla krótkiego ciągu dokładnego
$1\to N\xrightarrow{i}G\xrightarrow{p}Q\to1$
następujące warunki są równoważne:
\begin{enumerate}
\item istnieje przekrój homomorficzny $s:Q\to G$;
\item istnieje podgrupa $H\le G$ taka, że
$G=i(N)H$ i $i(N)\cap H=\{e\}$;
\item przy pewnym działaniu $\varphi:Q\to\operatorname{Aut}(N)$
istnieje izomorfizm $F:N\rtimes_\varphi Q\to G$
taki, że $F(n,e_Q)=i(n)$ dla każdego $n\in N$ oraz
$p(F(n,q))=q$ dla wszystkich $n\in N$, $q\in Q$.
\end{enumerate}
\end{twierdzenie}
\begin{proof}
Załóżmy (1) i połóżmy $H=s(Q)$. Ponieważ $s$ jest homomorfizmem,
$H$ jest podgrupą. Jeżeli $s(q)\in i(N)=\ker p$,
to $q=p(s(q))=e_Q$, zatem $s(q)=e_G$.
Mamy więc $i(N)\cap H=\{e\}$.
Dla dowolnego $g\in G$ element $s(p(g))$ ma ten sam obraz
w $Q$ co $g$; stąd $g s(p(g))^{-1}\in\ker p=i(N)$,
czyli $g\in i(N)H$. Zachodzi (2).

Zakładając (2), zauważamy, że ograniczenie $p|_H:H\to Q$
jest różnowartościowe, ponieważ jego jądro to
$H\cap\ker p=H\cap i(N)=\{e\}$.
Jest surjektywne: jeśli $q=p(g)$, zapiszmy $g=i(n)h$,
a wtedy $q=p(h)$. Zatem $p|_H$ jest izomorfizmem.
Jego odwrotność $s:Q\to H\subseteq G$ jest przekrojem
homomorficznym i wracamy do (1).

Przy (1) definiujemy $\varphi_q(n)$ jako jednoznaczny element
$N$, dla którego $i(\varphi_q(n))=s(q)i(n)s(q)^{-1}$.
Działanie jest dobrze określone, ponieważ $i(N)=\ker p$
jest normalne w $G$, a $i$ jest izomorfizmem na obraz.
Z Twierdzenia~\ref{aa:tw:internal-semidirect}
i podstawienia $h=s(q)$ otrzymujemy izomorfizm
$F(n,q)=i(n)s(q)$. Wprost mamy
$F(n,e_Q)=i(n)$ i $p(F(n,q))=q$, więc spełnia on (3).
Odwrotnie, izomorfizm z (3) przenosi standardowe włączenie
$q\mapsto(e_N,q)$ na przekrój $Q\to G$, więc (3) daje (1).
\end{proof}

\section{Przykłady: grupy diedralne, symetryczne i afiniczne}
\begin{przyklad}[Grupa diedralna jako iloczyn półprosty]
\label{aa:prz:dihedral-semidirect}
Niech $n\ge3$. Oznaczmy przez $r$ obrót foremnego $n$-kąta
o kąt $2\pi/n$, a przez $s$ odbicie w osi symetrii przechodzącej
przez jego środek. Zachodzą relacje
$r^n=e$, $s^2=e$ oraz $srs^{-1}=r^{-1}$.
Grupa $D_{2n}$ wszystkich symetrii ma $2n$ elementów:
$n$ obrotów $r^i$ oraz $n$ odbić $r^i s$.
Każdą symetrię można zapisać dokładnie w jednej z tych postaci;
obroty zachowują orientację płaszczyzny, a odbicia ją odwracają.
Podgrupa $N=\langle r\rangle\cong C_n$ jest normalna,
a $H=\langle s\rangle\cong C_2$ jest jej dopełnieniem,
gdyż $NH=D_{2n}$ i $N\cap H=\{e\}$.
Sprzężenie przez $s$ działa na $C_n$ przez inwersję:
$\varphi_s(r^i)=r^{-i}$.
Jest to automorfizm rzędu co najwyżej $2$, więc wyznacza
homomorfizm $C_2\to\operatorname{Aut}(C_n)$.
Twierdzenie~\ref{aa:tw:internal-semidirect} daje
$D_{2n}\cong C_n\rtimes_{\mathrm{inv}}C_2$.
Grupa nie jest na ogół iloczynem prostym, bo $srs^{-1}\ne r$
dla $n\ge3$.

Dla $n=4$ dostajemy symetrie kwadratu:
$D_8\cong C_4\rtimes C_2$, zgodnie z opisem
w Przykładzie~\ref{aa:prz:d8-nilpotent}.
Uwaga o notacji: stosujemy $D_{2n}$ na oznaczenie grupy
\emph{rzędu} $2n$; w części literatury tę samą grupę oznacza się $D_n$.
\end{przyklad}

\begin{przyklad}[Struktura $S_3$ i znaczenie normalności]
\label{aa:prz:s3-semidirect}
Niech $N=A_3=\langle r\rangle$,
$r=(1\,2\,3)$, oraz $H=\langle s\rangle$,
$s=(1\,2)$.
Podgrupa $N$ jest normalna, bo ma indeks dwa;
$|N|=3$, $|H|=2$, $N\cap H=\{e\}$.
Z Lematu~\ref{aa:lem:product-subgroups} otrzymujemy
$|NH|=3\cdot2=6=|S_3|$, więc $NH=S_3$.
Bezpośrednie obliczenie na punktach $1,2,3$ daje
$srs^{-1}=r^{-1}$.
Zatem $S_3\cong C_3\rtimes_{\mathrm{inv}}C_2$,
a nie $C_3\times C_2$: pierwsza grupa jest nieabelowa,
druga cykliczna rzędu $6$.
Nietrywialne działanie przez sprzężenie odróżnia tę grupę
od iloczynu prostego, mimo że w obu przypadkach rząd wynosi $6$.
\end{przyklad}

\begin{przyklad}[Grupa przekształceń afinicznych prostej nad ciałem]
\label{aa:prz:affine-group}
Niech $K$ będzie ciałem. Dla $a\in K$ i $u\in K^\times$
rozważmy przekształcenie $f_{a,u}:K\to K$,
$f_{a,u}(x)=ux+a$. Jest ono bijekcją, ponieważ
$u$ ma odwrotność; zachodzi
$f_{a,u}^{-1}(x)=u^{-1}(x-a)$.
Składając przekształcenia otrzymujemy
$f_{a,u}\circ f_{b,v}=f_{a+ub,uv}$.
Zatem grupa tych przekształceń jest iloczynem półprostym
$(K,+)\rtimes K^\times$, w którym
$u\in K^\times$ działa na $(K,+)$ przez mnożenie $x\mapsto ux$.

Dla $K=\mathbb F_7$ podgrupa
$\{f_{a,2^i}:a\in\mathbb F_7,\ i\in\{0,1,2\}\}$
ma $21$ elementów, ponieważ $2^3=1$ w $\mathbb F_7$.
Odzyskujemy więc dokładnie nieabelową grupę rzędu $21$
z Przykładu w Rozdziale~\ref{aa:ch:sylow}.
\end{przyklad}

\section{Przykład rozszerzenia, które się nie rozszczepia}
\begin{przyklad}[Rozszerzenie $C_2$ przez $C_2$ dające $C_4$]
\label{aa:prz:c4-nonsplit}
Rozważmy addytywny ciąg grup
\[
 0\longrightarrow C_2\xrightarrow{i}C_4
 \xrightarrow{p}C_2\longrightarrow0,
 \qquad i([a]_2)=[2a]_4,\quad p([b]_4)=[b]_2.
\]
Wszystkie odwzorowania są homomorfizmami:
$i$ jest iniektywne, $p$ surjektywne,
a $\ker p=\{[0]_4,[2]_4\}=\operatorname{im}i$.
Ciąg jest więc dokładny.

Załóżmy, że istnieje przekrój homomorficzny
$s:C_2\to C_4$ z $p\circ s=\operatorname{id}_{C_2}$.
Wówczas $p(s([1]_2))=[1]_2$, czyli $s([1]_2)$
musi być jedną z klas $[1]_4$ lub $[3]_4$.
Obie mają rząd $4$. Z drugiej strony $[1]_2$ ma rząd $2$,
a homomorfizm musi spełniać
$2s([1]_2)=s(2[1]_2)=s([0]_2)=[0]_4$,
co wyklucza element rzędu $4$. Otrzymujemy sprzeczność.
Ciąg nie rozszczepia się, choć wszystkie trzy grupy są abelowe.

Dla porównania $C_2\times C_2$ daje rozszczepiony ciąg
$0\to C_2\to C_2\times C_2\to C_2\to0$,
gdzie włączenie i rzut dotyczą odpowiednio pierwszego
i drugiego czynnika, a przekrojem jest $[b]\mapsto(0,[b])$.
Mamy więc dwa nieizomorficzne rozszerzenia tych samych grup.
\end{przyklad}

\begin{przyklad}[Nierozszczepiające się rozszerzenie z nieabelową grupą środkową]
\label{aa:prz:q8-nonsplit}
Otrzymamy inny przykład, w którym iloraz i jądro są abelowe,
ale cała grupa nie jest abelowa. Niech $E$ będzie macierzą
jednostkową stopnia $2$, $\mathrm i\in\mathbb C$ jednostką
urojoną, a
\[
 I=\begin{pmatrix}\mathrm i&0\\0&-\mathrm i\end{pmatrix},
 \qquad
 J=\begin{pmatrix}0&1\\-1&0\end{pmatrix}.
\]
Bezpośrednie mnożenie macierzy daje
$I^2=J^2=-E$ i $JI=-IJ$. Z tego wynika, że
$Q_8=\{\pm E,\pm I,\pm J,\pm IJ\}$
jest grupą z mnożeniem macierzy: iloczyn dowolnych dwóch
elementów, przy użyciu podanych relacji, sprowadza się
do jednego z ośmiu wypisanych; odwrotności wynikają
z $I^{-1}=-I$, $J^{-1}=-J$.
Osiem macierzy jest różnych, zatem $|Q_8|=8$.

Podgrupa $Z=\{E,-E\}$ jest centralna i izomorficzna z $C_2$.
Klasy $IZ,JZ$ w $Q_8/Z$ mają rząd $2$, ponieważ
$I^2=J^2=-E\in Z$. Komutują, gdyż $JI=-IJ$, a ich
klasy są różne i nie są neutralne: $I,J\notin Z$ i
$I\notin JZ$. Ponadto $(IZ)(JZ)=IJZ$ jest różne
od $Z$, $IZ$ i $JZ$, ponieważ osiem wypisanych macierzy
tworzy cztery rozłączne dwuelementowe warstwy modulo $Z$.
Te cztery klasy wyczerpują $Q_8/Z$, zatem iloraz
jest generowany przez dwa komutujące elementy rzędu $2$, skąd
$Q_8/Z\cong C_2\times C_2$. Powstaje krótki ciąg dokładny
\[
 1\longrightarrow Z\longrightarrow Q_8
 \longrightarrow Q_8/Z\longrightarrow1.
\]
Ten ciąg nie rozszczepia się. Istotnie, przekrój homomorficzny
$s:Q_8/Z\to Q_8$ byłby iniektywny, ponieważ
$p\circ s=\operatorname{id}_{Q_8/Z}$. Musiałby więc
umieszczać w $Q_8$ podgrupę izomorficzną z
$C_2\times C_2$, która ma trzy różne elementy rzędu $2$.
Tymczasem jedynym elementem rzędu $2$ w $Q_8$ jest $-E$:
dla każdego $X\in\{\pm I,\pm J,\pm IJ\}$ mamy
$X^2=-E$, a więc $X$ ma rząd $4$.
To sprzeczność. Rozszczepienie jest więc własnością
konkretnego ciągu dokładnego, nie tylko rzędów jego grup.
\end{przyklad}

\begin{uwaga}[Słowo \emph{rozszczepienie} jest ważniejsze niż sam rozmiar]
Każda surjekcja zbiorów skończonych ma przekrój
\emph{mnogościowy}, ale nie każdy taki przekrój jest homomorfizmem.
Dla iloczynu półprostego potrzebujemy przekroju zachowującego
mnożenie. Przykład~\ref{aa:prz:c4-nonsplit} pokazuje,
że tego dodatkowego wymagania nie wolno pomijać.
\end{uwaga}

\section{Grupy rzędu iloczynu dwóch różnych liczb pierwszych}
\begin{lemat}[Automorfizmy grupy cyklicznej rzędu pierwszego]
\label{aa:lem:aut-cq}
Jeżeli $q$ jest liczbą pierwszą, to
$\operatorname{Aut}(C_q)\cong(\mathbb Z/q\mathbb Z)^\times$.
W szczególności $|\operatorname{Aut}(C_q)|=q-1$.
Dla każdej liczby pierwszej $p$ dzielącej $q-1$
grupa $\operatorname{Aut}(C_q)$ ma dokładnie jedną
podgrupę rzędu $p$.
\end{lemat}
\begin{proof}
Grupa $C_q$ jest generowana przez $[1]_q$.
Homomorfizm $f:C_q\to C_q$ jest zatem wyznaczony przez
$f([1]_q)=[a]_q$ i ma postać $f([k]_q)=[ak]_q$.
Jest automorfizmem dokładnie wtedy, gdy $[a]_q\ne[0]_q$,
ponieważ dla $q$ pierwszego każda niezerowa klasa jest
odwracalna. Składanie automorfizmów odpowiada mnożeniu klas,
co daje izomorfizm i liczbę $q-1$.

Niech $p\mid(q-1)$. Twierdzenie Cauchy'ego zapewnia istnienie
elementu rzędu $p$, a więc podgrupy rzędu $p$.
Jej wszystkie elementy są rozwiązaniami równania $x^p=1$
w ciele $\mathbb F_q$. Wielomian $X^p-1$ ma nad ciałem
co najwyżej $p$ pierwiastków: dla dowolnego pierwiastka $a$
można wyłączyć czynnik $(X-a)$, a następnie stosować ten
argument indukcyjnie do wielomianu o mniejszym stopniu.
Podgrupa rzędu $p$ zawiera już $p$ różnych pierwiastków
tego wielomianu, więc jest całym zbiorem rozwiązań $x^p=1$.
Każda inna podgrupa rzędu $p$ musiałaby być zawarta w tym
samym zbiorze i miałaby $p$ elementów, więc byłaby równa pierwszej.
\end{proof}

\begin{twierdzenie}[Klasyfikacja grup rzędu $pq$]
\label{aa:tw:classify-pq}
Niech $p<q$ będą różnymi liczbami pierwszymi.
\begin{enumerate}
\item Jeśli $p\nmid(q-1)$, to z dokładnością do izomorfizmu
istnieje tylko jedna grupa rzędu $pq$, mianowicie $C_{pq}$.
\item Jeśli $p\mid(q-1)$, to z dokładnością do izomorfizmu
istnieją dokładnie dwie grupy rzędu $pq$:
cykliczna $C_{pq}$ i nieabelowa $C_q\rtimes C_p$,
w której generator $C_p$ działa na $C_q$
przez automorfizm rzędu $p$.
\end{enumerate}
\end{twierdzenie}
\begin{proof}
Niech $G$ będzie grupą rzędu $pq$.
Z trzeciego twierdzenia Sylowa $n_q\mid p$ i
$n_q\equiv1\pmod q$. Ponieważ $p<q$, jedyna możliwość
to $n_q=1$. Podgrupa $Q$ rzędu $q$ jest więc normalna.
Z pierwszego twierdzenia Sylowa istnieje podgrupa $P$ rzędu $p$.
Obie są cykliczne, a $P\cap Q=\{e\}$, bo rząd przekroju
dzieli jednocześnie $p$ i $q$.
Lemat~\ref{aa:lem:product-subgroups} daje
$|QP|=|Q|\,|P|=pq=|G|$, więc $QP=G$.
Z Twierdzenia~\ref{aa:tw:internal-semidirect}
$G\cong C_q\rtimes_\varphi C_p$ dla pewnego homomorfizmu
$\varphi:C_p\to\operatorname{Aut}(C_q)$.

Ponieważ $C_p$ jest generowana przez element rzędu $p$,
obraz $\varphi$ ma rząd $1$ albo $p$.
Pierwszy przypadek oznacza działanie trywialne,
więc $G\cong C_q\times C_p\cong C_{pq}$;
ostatni izomorfizm wynika z Przykładu o $C_6$,
a jego dowód działa bez zmiany dla różnych liczb pierwszych.
Jeśli obraz ma rząd $p$, to twierdzenie Lagrange'a zastosowane
do $\operatorname{Aut}(C_q)$ daje $p\mid q-1$.
Stąd przy $p\nmid(q-1)$ istnieje tylko grupa cykliczna.

Załóżmy $p\mid(q-1)$. Z Lematu~\ref{aa:lem:aut-cq}
istnieje automorfizm $\alpha$ rzędu $p$.
Wysłanie generatora $C_p$ na $\alpha$ określa nietrywialne
działanie, a więc grupę $C_q\rtimes_\alpha C_p$ rzędu $pq$.
Grupa ta nie jest abelowa: wybierzmy $x\in C_q$ taki,
że $\alpha(x)\ne x$ (istnieje, bo $\alpha$ nie jest identycznością).
Wtedy $(e_{C_q},1)(x,0)(e_{C_q},1)^{-1}=(\alpha(x),0)\ne(x,0)$.

Pokażemy, że wszystkie nietrywialne działania dają grupy
izomorficzne. Z Lematu~\ref{aa:lem:aut-cq} każdy ich obraz
jest tą samą podgrupą $U\le\operatorname{Aut}(C_q)$ rzędu $p$.
Dwa homomorfizmy $\varphi,\psi:C_p\to U$ o obrazach równych $U$
są izomorfizmami na $U$.
Zatem $\theta=\varphi^{-1}\circ\psi:C_p\to C_p$ jest
automorfizmem i $\psi(h)=\varphi(\theta(h))$.
Wzór $(x,h)\mapsto(x,\theta(h))$ jest wtedy izomorfizmem
$C_q\rtimes_\psi C_p\to C_q\rtimes_\varphi C_p$:
po obu stronach pierwszy czynnik iloczynu jest
$x\,\psi_h(y)=x\,\varphi_{\theta(h)}(y)$,
a drugi jest zgodny, bo $\theta$ jest homomorfizmem.
Izomorfizm jest bijektywny, gdyż $\theta$ jest automorfizmem.
To kończy klasyfikację.
\end{proof}

\begin{przyklad}[Wszystkie grupy rzędów $15$ i $21$]
Dla $|G|=15=3\cdot5$ mamy $3\nmid(5-1)=4$,
więc $G\cong C_{15}$: nie ma nieabelowej grupy rzędu $15$.
Dla $|G|=21=3\cdot7$ zachodzi $3\mid(7-1)$,
więc są dokładnie dwie grupy: $C_{21}$ i
$C_7\rtimes C_3$. W drugiej generator $C_3$ działa
na $C_7$ przez mnożenie przez $2$, ponieważ
$2^3\equiv1\pmod7$ i $2\not\equiv1\pmod7$.
To dokładnie grupa przekształceń $f_{a,i}(x)=2^i x+a$
z poprzedniego przykładu~\ref{aa:prz:affine-group}.
Dla pełnego rachunku możemy zapisać jej elementy jako
pary $(a,i)\in C_7\times C_3$ z mnożeniem
\[
 (a,i)(b,j)=(a+2^i b,\ i+j),
\]
gdzie pierwszą współrzędną redukujemy modulo $7$, drugą
modulo $3$. Wzór jest dobrze określony: zastąpienie
$i$ przez $i+3$ nie zmienia mnożnika $2^i$ modulo $7$,
bo $2^3\equiv1\pmod 7$. Elementami $r=(1,0)$
i $s=(0,1)$ generujemy całą grupę: potęgi $r$ dają
wszystkie $(a,0)$, a potęgi $s$ dają $(0,i)$.
Obliczamy $srs^{-1}=(2,0)=r^2$, a zatem
$sr\ne rs$. Przykład pokazuje zarówno sposób kodowania
rozszerzenia przez działanie, jak i źródło jego nieprzemienności.
\end{przyklad}

\section{Skończone grupy nilpotentne i iloczyny podgrup Sylowa}
\begin{lemat}[Warunek normalizatora w grupie nilpotentnej]
\label{aa:lem:nilpotent-normalizer}
Niech $G$ będzie grupą nilpotentną i $H<G$ właściwą
podgrupą. Wtedy $H$ jest właściwą podgrupą swojego
normalizatora: $H<N_G(H)$.
\end{lemat}
\begin{proof}
Ponieważ $G$ jest nilpotentna, istnieje $c$ z $Z_c(G)=G$
na mocy Lematu~\ref{aa:lem:upper-lower}.
Wybierzmy najmniejsze $i\ge1$, dla którego $Z_i(G)$
nie zawiera się w $H$. Taki indeks istnieje, ponieważ
$Z_0(G)=\{e\}\subseteq H$, natomiast $Z_c(G)=G\nsubseteq H$.
Wówczas $Z_{i-1}(G)\subseteq H$.
Wybierzmy $z\in Z_i(G)\setminus H$.
Z definicji $Z_i(G)$ dla każdego $h\in H$ zachodzi
$[z,h]\in Z_{i-1}(G)\subseteq H$.
Stąd $zhz^{-1}=[z,h]h\in H$.
Ten sam argument stosujemy do $z^{-1}\in Z_i(G)$:
$z^{-1}hz\in H$ dla wszystkich $h\in H$.
Zatem $zHz^{-1}\subseteq H$ i $z^{-1}Hz\subseteq H$.
Sprzężenie pierwszej inkluzji przez $z^{-1}$ daje
$H\subseteq z^{-1}Hz$, więc w rzeczywistości
$z^{-1}Hz=H$, a tym samym $zHz^{-1}=H$.
Wynika stąd $z\in N_G(H)\setminus H$.
Każdy element $H$ normalizuje $H$, zatem
$H\subsetneq N_G(H)$.
\end{proof}

\begin{twierdzenie}[Charakterystyka skończonych grup nilpotentnych]
\label{aa:tw:nilpotent-sylow-product}
Dla skończonej grupy $G$ następujące warunki są równoważne:
\begin{enumerate}
\item $G$ jest nilpotentna;
\item każda podgrupa Sylowa grupy $G$ jest normalna;
\item $G$ jest izomorficzna z iloczynem prostym swoich podgrup
Sylowa, po jednej dla każdego dzielnika pierwszego $|G|$.
\end{enumerate}
\end{twierdzenie}
\begin{proof}
Jeżeli $G=\{e\}$, wszystkie warunki są spełnione, a pusty
iloczyn prosty interpretujemy jako grupę trywialną.
W dalszym dowodzie zakładamy $|G|>1$.

\emph{(1) $\Rightarrow$ (2).}
Niech $P$ będzie podgrupą $p$-Sylowa i
$H=N_G(P)$ jej normalizatorem.
Mamy $P\trianglelefteq H$ z definicji normalizatora.
Ponieważ $P\le H\le G$, największa potęga $p$
dzieląca $|H|$ jest taka sama jak dla $|G|$.
Zatem $P$ jest podgrupą $p$-Sylowa także w $H$.
Przypuśćmy $H\ne G$.
Z poprzedniego lematu istnieje element
$x\in N_G(H)\setminus H$.
Ponieważ $x$ normalizuje $H$, podgrupa $xPx^{-1}$
zawiera się w $H$ i ma ten sam rząd co $P$;
jest więc podgrupą $p$-Sylowa $H$.
Ale $P$ jest normalna w $H$, a wszystkie podgrupy $p$-Sylowa
w $H$ są sprzężone (drugie twierdzenie Sylowa), więc
$P$ jest jedyną podgrupą $p$-Sylowa $H$.
Musi więc zachodzić $xPx^{-1}=P$, czyli $x\in N_G(P)=H$,
wbrew wyborowi $x$. Zatem $H=G$, więc $P\trianglelefteq G$.
Ponieważ $p$ i $P$ były dowolne, wszystkie podgrupy Sylowa
są normalne.

\emph{(2) $\Rightarrow$ (3).}
Zapiszmy $|G|=p_1^{a_1}\cdots p_r^{a_r}$ dla parami różnych
liczb pierwszych i wybierzmy podgrupy $P_i$ rzędu $p_i^{a_i}$.
Wszystkie są normalne z założenia.
Dla $i\ne j$ mamy $P_i\cap P_j=\{e\}$,
gdyż rząd przecięcia dzieli obie względnie pierwsze potęgi.
Dla $x\in P_i$ i $y\in P_j$ komutator $[x,y]$ należy
zarówno do $P_i$, jak i do $P_j$, ponieważ obie podgrupy
są normalne. Zatem $[x,y]\in P_i\cap P_j=\{e\}$,
czyli elementy $P_i$ i $P_j$ komutują.

Indukcyjnie zbudujemy podgrupę $K_j=P_1\cdots P_j$
o rzędzie $p_1^{a_1}\cdots p_j^{a_j}$, izomorficzną
z iloczynem prostym $P_1\times\cdots\times P_j$.
Dla $j=1$ jest to oczywiste.
Jeśli własność zachodzi dla $j$, to $K_j$ jest normalna,
gdyż iloczyn normalnych podgrup pozostaje normalny.
Rzędy $K_j$ oraz $P_{j+1}$ są względnie pierwsze,
więc $K_j\cap P_{j+1}=\{e\}$.
Twierdzenie~\ref{aa:tw:internal-direct} daje
$K_jP_{j+1}\cong K_j\times P_{j+1}$
i rząd równy iloczynowi rzędów.
Po $r$ krokach mamy podgrupę $K_r$ o rzędzie
$\prod_i p_i^{a_i}=|G|$, więc $K_r=G$.
To daje (3).

\emph{(3) $\Rightarrow$ (1).}
Każda $P_i$ jest skończoną $p_i$-grupą,
więc jest nilpotentna na mocy
Twierdzenia~\ref{aa:tw:pgroup-nilpotent}.
Centrum iloczynu prostego to iloczyn centrów czynników
(Twierdzenie~\ref{aa:tw:direct-product}).
Indukcja po $j$ i definicja górnego szeregu centralnego
dają
$Z_j(P_1\times\cdots\times P_r)
=Z_j(P_1)\times\cdots\times Z_j(P_r)$:
dla $j=0$ jest to równość grup trywialnych;
przy kroku indukcyjnym iloraz iloczynu przez $Z_j$
identyfikuje się z iloczynem ilorazów $P_i/Z_j(P_i)$,
a centrum tego iloczynu jest iloczynem centrów.
Dla każdego $i$ istnieje $c_i$ takie, że $Z_{c_i}(P_i)=P_i$.
Biorąc $c=\max_i c_i$, otrzymujemy $Z_c(G)=G$,
a stąd nilpotentność na mocy
Lematu~\ref{aa:lem:upper-lower}.
\end{proof}

\begin{przyklad}[Rozwiązalność bez nilpotentności jeszcze raz]
Grupa $S_3$ ma podgrupy Sylowa rzędów $3$ i $2$.
Podgrupa rzędu $3$ jest normalna, lecz trzy podgrupy
rzędu $2$ są wzajemnie sprzężone i żadna nie jest normalna.
Na mocy poprzedniego twierdzenia $S_3$ nie jest nilpotentna,
choć jest rozwiązalna. To daje drugi, niezależny od centrum,
sposób rozpoznania różnicy pomiędzy tymi dwiema własnościami.
\end{przyklad}

\section{Związek z iloczynem tensorowym grup abelowych}
\label{aa:sec:tensor-group-plan}
Dla \emph{grup abelowych} $A,B$ istnieje naturalna konstrukcja
$A\otimes_{\mathbb Z}B$, ale nie jest ona ani iloczynem prostym,
ani opisanym wyżej iloczynem półprostym.
Podstawą jest obserwacja, że każda grupa abelowa jest
kanonicznie modułem nad $\mathbb Z$: definiujemy
$n\cdot a$ jako $n$-krotną sumę $a$ dla $n>0$,
$0\cdot a=0$ oraz $(-n)\cdot a=-(n\cdot a)$.
Wówczas iloczyn tensorowy grup abelowych jest po prostu
iloczynem tensorowym $\mathbb Z$-modułów.
Jego własność uniwersalna opisuje odwzorowania
addytywne w każdej zmiennej osobno.

Dla dodatnich liczb całkowitych $m,n$ mamy
$C_m\otimes_{\mathbb Z}C_n\cong C_{\gcd(m,n)}$.
Przykładowo $C_2\otimes_{\mathbb Z}C_3=0$, mimo że
$C_2\times C_3\cong C_6$. Natomiast
$C_4\otimes_{\mathbb Z}C_6\cong C_2$, choć
$C_4\times C_6$ ma $24$ elementy.

W tym skrypcie konstrukcję przez generatory i relacje oraz
dowód powyższego wzoru przedstawiamy już w
Rozdziale~\ref{aa:ch:grupy-abelowe},
w szczególności w
Twierdzeniu~\ref{aa:tw:abelian-tensor-cyclic}.
Ogólną konstrukcję dla modułów nad pierścieniem przemiennym
oraz jej własności rozwijamy dalej w
Rozdziale~\ref{aa:ch:module-tensor-product}.
Odsyłacze te prowadzą do rzeczywistych rozdziałów tego
samego skryptu; nie są zależnością od obcego modułu ani
od zewnętrznego pliku pomocniczego.
